% ============================================================================
% Fichier  : partie0/P00_une_nuit_a_garoua.tex
% Titre    : Prologue : Une Nuit à Garoua
% Auteur   : MINKA MI NGUIDJOÏ Thierry Emmanuel
% Version  : 1.0 -- novembre 2026
% ============================================================================
% Opération GAROUA : réalisation du scénario d'attaque décrit dans INF3226,
% Chapitre 5 (hameçonnage du prestataire de Garoua, accès distant sans
% double authentification, serveur non corrigé, sauvegardes en ligne).
% Personnages communs au corpus : M. ETEME (DG), Mme NKENG (contrôleuse
% interne), M. KOUNOU (responsable informatique).
% ============================================================================

\chapter{Prologue : Une Nuit à Garoua}
\label{chap:prologue}
\chaptermark{Une nuit à Garoua}

\epigraph{%
    « On ne protège bien que ce que l'on connaît,\\
    et l'on ne connaît bien que ce que l'on a mesuré. »%
}{Principe du cours}

\niveauChapitre{Introductif}{Aucun prérequis~: ce prologue présente le cas
qui sert de fil rouge au cours et les questions auxquelles il apprend à
répondre.}

\begin{boxobjectifs}
\begin{itemize}
    \item Suivre, heure par heure, le déroulement d'une attaque par
          rançongiciel dans une entreprise camerounaise.
    \item Repérer, dans ce récit, les défaillances qui l'ont rendue possible.
    \item Distinguer une constatation, une hypothèse et une opinion.
    \item Situer chacune des questions soulevées dans une partie du cours.
    \item Identifier les métiers de la sécurité et leur rôle dans la crise.
\end{itemize}
\end{boxobjectifs}

\bigskip

% ============================================================================
\section*{Question directrice : comment une entreprise perd-elle tout en une nuit~?}
% ============================================================================

LiZbiZ\footnote{LiZbiZ, ses dirigeants et ses prestataires sont fictifs~;
les dates et les montants aussi. Les mécanismes d'attaque décrits sont
exacts et documentés dans les sources citées au fil du cours.} est une
entreprise de logistique de Douala. Elle stocke, achemine et livre des
marchandises pour des importateurs et des distributeurs, dans douze agences.
Tout repose sur une application intégrée, hébergée dans quatre centres de
données (Douala, Yaoundé, Garoua, Bafoussam). Un vendredi de mars, à
21~h~40, la facturation s'arrête. Le lundi suivant, l'entreprise ne peut
plus ni facturer ni préparer une expédition.

\medskip
\textbf{Plan du prologue.}
\begin{enumerate}
    \item LiZbiZ, son système d'information et ses acteurs (section~\ref{sec:p00-lizbiz}).
    \item La nuit du 5 au 6 mars (section~\ref{sec:p00-nuit}).
    \item Ce que l'on sait le lundi matin (section~\ref{sec:p00-lundi}).
    \item Trente questions, six parties (section~\ref{sec:p00-questions}).
    \item Les métiers de la sécurité (section~\ref{sec:p00-metiers}).
\end{enumerate}

\separateur

% ============================================================================
\section{LiZbiZ, son Système d'Information et ses Acteurs}
\label{sec:p00-lizbiz}
% ============================================================================

\begin{tableaucours}{L{3.6cm}L{8.6cm}}{%
    \tableauHeader{Élément} &
    \tableauHeader{Description}
}
Activité & Logistique~: entreposage, transport, livraison \\
Implantation & Siège à Douala~; douze agences~; environ 150~salariés \\
Informatique & Quatre centres de données~; une application intégrée
(stocks, expéditions, facturation, paiements)~; messagerie en nuage \\
Direction & M.~ETEME, directeur général~; Mme~NKENG, contrôleuse
interne~; M.~KOUNOU, responsable informatique \\
Sécurité & Aucun responsable de la sécurité désigné~; un pare-feu par
centre~; un antivirus sur les postes \\
Prestataire & Depuis septembre~2026, la société SAHEL INFOGÉRANCE exploite
les centres de Garoua et de Bafoussam, avec un accès distant \\
\end{tableaucours}
\vspace{-0.8em}
\begin{center}{\small\itshape Tableau~P.1. Fiche d'identité de LiZbiZ}\end{center}

L'organisation de LiZbiZ ressemble à celle de nombreuses entreprises de sa
taille~: une informatique qui a grandi avec l'activité, un responsable
informatique qui fait tout, des prestataires appelés au fil des besoins, et
une sécurité confiée aux produits plutôt qu'à des personnes. Le cours ne
juge pas cette organisation~; il apprend à la décrire, à mesurer ses
risques et à la faire évoluer.

\begin{boxinfo}
\textbf{Un cas, un corpus.} LiZbiZ est aussi le cas principal du cours
d'audit des systèmes d'information (UE INF3226), qui en audite la
gouvernance et les contrôles. Ici, on la regarde avec les yeux de celui
qui doit la protéger.
\end{boxinfo}

% ============================================================================
\section{La Nuit du 5 au 6 Mars}
\label{sec:p00-nuit}
% ============================================================================

La chronologie ci-dessous a été reconstituée après coup, à partir des
journaux qui avaient survécu. Les heures sont celles de Douala.

\begin{tableaucours}{L{2.4cm}L{9.8cm}}{%
    \tableauHeader{Date et heure} &
    \tableauHeader{Événement (constaté après coup)}
}
19 février & Un technicien de SAHEL INFOGÉRANCE reçoit un courriel qui
imite une alerte de son fournisseur de messagerie, et saisit son mot de
passe sur une fausse page. \\
2 mars, 23~h~12 & Connexion à la passerelle d'accès distant de Garoua avec
le compte du technicien~; aucune double authentification n'est exigée. \\
3 et 4 mars & Exploration du réseau interne~: découverte d'un serveur de
fichiers dont un correctif critique, publié huit mois plus tôt, n'a pas
été installé. \\
4 mars, 02~h~05 & Obtention de droits d'administrateur sur ce serveur,
puis sur le domaine. \\
5 mars, 20~h~30 & Copie de dossiers clients vers un service de stockage
externe (environ 40~Go). \\
5 mars, 21~h~15 & Suppression des sauvegardes, stockées sur un partage
réseau accessible depuis le domaine. \\
5 mars, 21~h~40 & Chiffrement des serveurs de Garoua et de Douala~; un
message de rançon s'affiche~: 85~000~dollars en cryptoactifs, sous peine
de publier les données des clients. \\
5 mars, 22~h~05 & Le technicien d'astreinte, M.~ABENA, constate que
l'application ne répond plus et redémarre les serveurs. \\
6 mars, 07~h~30 & M.~KOUNOU découvre l'ampleur de l'attaque et prévient le
directeur général. \\
\end{tableaucours}
\vspace{-0.8em}
\begin{center}{\small\itshape Tableau~P.2. Chronologie de l'Opération GAROUA}\end{center}

\begin{boxalert}
\textbf{Violation typique.} Redémarrer les serveurs « pour voir si cela
repart », comme M.~ABENA à 22~h~05.

\smallskip
\textbf{Règle~:} un système compromis s'isole du réseau, il ne se redémarre
pas. Le redémarrage efface la mémoire vive, et avec elle des traces
précieuses (processus, connexions, parfois des clés). Le Chapitre~\ref{chap:operations}
et le Chapitre~\ref{chap:investigation} disent quoi faire à la place.
\end{boxalert}

% ============================================================================
\section{Ce que l'on Sait le Lundi Matin}
\label{sec:p00-lundi}
% ============================================================================

Le lundi 8 mars, à 8~h, M.~ETEME réunit son comité de direction. Chacun
parle avec assurance. Il faut pourtant trier ce qui est dit.

\begin{tableaucours}{L{6.4cm}L{2.4cm}L{3.4cm}}{%
    \tableauHeader{Affirmation} &
    \tableauHeader{Nature} &
    \tableauHeader{Ce qu'il faudrait}
}
« Les serveurs de Garoua et de Douala sont chiffrés. » & Constatation &
Rien de plus~: c'est visible \\
« L'attaque vient du prestataire. » & Hypothèse & Les journaux de la
passerelle d'accès distant \\
« Les données des clients ont été volées. » & Hypothèse & Les journaux du
pare-feu, le volume sortant \\
« On peut tout restaurer depuis les sauvegardes. » & Opinion, fausse &
Un test de restauration~: les sauvegardes ont été supprimées \\
« Personne n'aurait pu l'empêcher. » & Opinion & Une analyse des mesures
manquantes \\
\end{tableaucours}
\vspace{-0.8em}
\begin{center}{\small\itshape Tableau~P.3. Trier les affirmations du comité}\end{center}

\begin{boxdef}
\termecle{Constatation}~: fait vérifiable par un tiers, appuyé sur une
trace conservée. \termecle{Hypothèse}~: explication plausible, qui reste à
vérifier. \termecle{Opinion}~: jugement qui ne repose ni sur une trace ni
sur une démarche.
\end{boxdef}

Le cours demande de toujours séparer ces trois registres. Un responsable
de la sécurité qui présente une hypothèse comme une constatation conduit
sa direction à de mauvaises décisions~; un responsable qui ne dit pas ce
qu'il ignore lui retire le moyen de décider.

\begin{boxscenario}{Opération GAROUA. Les trois décisions du lundi}
\textbf{Situation.} Le comité doit décider avant midi~: payer ou non la
rançon~; prévenir ou non les clients et les autorités~; reconstruire ou
restaurer.

\medskip
\textbf{Ce qui manque.} Personne ne sait quelles données sont sorties, ni
depuis quand l'attaquant était présent, ni si d'autres centres sont
touchés. LiZbiZ n'a ni plan de réponse aux incidents, ni plan de
continuité, ni contrat qui oblige SAHEL INFOGÉRANCE à fournir ses journaux.
Elle doit décider dans le brouillard.
\end{boxscenario}

% ============================================================================
\section{Trente Questions, Six Parties}
\label{sec:p00-questions}
% ============================================================================

L'Opération GAROUA soulève des questions qu'aucun participant au comité ne
sait encore traiter. Le cours y répond partie par partie.

\begin{tableaucours}{L{2.4cm}L{7.6cm}L{2.2cm}}{%
    \tableauHeader{Partie} &
    \tableauHeader{Questions soulevées par l'Opération GAROUA} &
    \tableauHeader{Chapitres}
}
I. Fondements & Qu'est-ce qui a été atteint~: la confidentialité,
l'intégrité, la disponibilité~? Qu'est-ce qu'un rançongiciel, une
élévation de privilèges~? Le chiffrement aurait-il protégé les données
volées~? & \ref{chap:proprietes} à~\ref{chap:attaques} \\
II. Gouvernance et risques & Qui était responsable de la sécurité~? Le
risque avait-il été apprécié~? Quelles données étaient les plus sensibles~?
Le technicien avait-il été formé~? & \ref{chap:gouvernance} à~\ref{chap:humain} \\
III. Cadre normatif et juridique & Quelles mesures une norme aurait-elle
exigées~? Que dit la loi camerounaise sur la notification, sur la
conservation des journaux~? Peut-on payer une rançon~? & \ref{chap:normes}
à~\ref{chap:ethique} \\
IV. Sécurité technique & Pourquoi l'accès distant était-il à un seul
facteur~? Pourquoi le réseau n'était-il pas segmenté~? Pourquoi le
correctif n'était-il pas installé~? Pourquoi personne n'a vu la copie de
40~Go~? & \ref{chap:modeles} à~\ref{chap:operations} \\
V. Défis avancés & Comment reprendre l'activité~? Comment préserver les
preuves pour une plainte~? Une IA aurait-elle détecté l'attaque~? &
\ref{chap:continuite} à~\ref{chap:emergents} \\
VI. Mesurer & Comment savoir, demain, si LiZbiZ est mieux protégée~? &
\ref{chap:mesurer} \\
\end{tableaucours}
\vspace{-0.8em}
\begin{center}{\small\itshape Tableau~P.4. Les questions de l'Opération GAROUA et
les parties du cours}\end{center}

\begin{boxinfo}
\textbf{Ce que dit la recherche.} Une organisation qui ne sait pas
répondre à une question sur ses propres systèmes ne doit pas faire comme
si la réponse était favorable. Le corpus de l'auteur appelle cette
réponse \NEK{} (« connaissance insuffisante »), à côté de \OUI{} et de
\NON{}~\cite{minka2026nek}. Le lundi matin, la plupart des réponses de
LiZbiZ sont des \NEK{}~: c'est le point de départ, non une fatalité.
\end{boxinfo}

% ============================================================================
\section{Les Métiers de la Sécurité}
\label{sec:p00-metiers}
% ============================================================================

Cinq métiers reviennent tout au long du cours. Chacun aurait pu, à sa
place, empêcher ou limiter l'Opération GAROUA.

\begin{boxmetier}{RSSI}
\textbf{Le responsable de la sécurité des systèmes d'information} aurait
apprécié le risque lié au prestataire, exigé la double authentification et
une clause d'accès aux journaux, et fait tester les sauvegardes. LiZbiZ
n'en avait pas.
\end{boxmetier}

\begin{boxmetier}{Administrateur systèmes et réseaux}
\textbf{L'administrateur} aurait installé le correctif critique, séparé les
sauvegardes du domaine et segmenté le réseau entre les centres.
\end{boxmetier}

\begin{boxmetier}{Analyste SOC}
\textbf{L'analyste d'un centre opérationnel de sécurité} aurait vu la
connexion de 23~h~12, l'élévation de privilèges et la copie de 40~Go, à
condition que les journaux soient collectés et surveillés.
\end{boxmetier}

\begin{boxmetier}{Développeur}
\textbf{Le développeur} de l'application intégrée aurait prévu une
authentification forte, des journaux exploitables et une reprise
possible à partir de données saines.
\end{boxmetier}

\begin{boxmetier}{Auditeur de sécurité}
\textbf{L'auditeur de sécurité} aurait relevé, avant l'attaque,
l'absence de double authentification, le correctif manquant et la
fragilité des sauvegardes, et l'aurait écrit dans un rapport adressé à la
direction.
\end{boxmetier}

\separateur

\begin{boxresume}
\textbf{Ce qu'il faut retenir de ce prologue~:}
\begin{itemize}
    \item \textbf{Une attaque est une chaîne}~: hameçonnage, accès distant,
          exploration, élévation de privilèges, exfiltration, destruction
          des sauvegardes, chiffrement. Chaque maillon était une occasion de
          l'arrêter.
    \item \textbf{Constatation, hypothèse, opinion} se distinguent~; une
          décision prise sur une opinion présentée comme un fait est une
          mauvaise décision.
    \item \textbf{Ne pas savoir} est une réponse à signaler (\NEK{}), pas à
          masquer.
    \item \textbf{La sécurité est un métier}, réparti entre plusieurs
          fonctions~; un produit ne remplace aucune d'elles.
\end{itemize}
\end{boxresume}

\begin{boxinfo}
\textbf{Points de contrôle de l'élève.} Je sais~: raconter l'Opération
GAROUA en sept étapes~; classer une affirmation en constatation, hypothèse
ou opinion~; associer une question de sécurité à une partie du cours~;
dire ce que chacun des cinq métiers aurait pu faire.
\end{boxinfo}

\bigskip

\begin{boxexo}[title={\ding{45}\quad Exercice~P.1. La chaîne d'attaque \niveaubase}]
Recopiez les sept étapes de l'attaque (Tableau~P.2) et indiquez, pour
chacune, une mesure qui l'aurait empêchée ou détectée.
\end{boxexo}

\begin{boxexo}[title={\ding{45}\quad Exercice~P.2. Trier les affirmations \niveaubase}]
Classez en constatation, hypothèse ou opinion~: (a) « l'attaquant est
russe »~; (b) « le message de rançon demande 85~000 dollars »~; (c) « le
prestataire est incompétent »~; (d) « le compte du technicien a servi à
23~h~12 »~; (e) « les données des clients seront publiées ».
\end{boxexo}

\begin{boxexo}[title={\ding{45}\quad Exercice~P.3. Ce qui a été atteint \niveaubase}]
Pour chacune des propriétés confidentialité, intégrité et disponibilité,
dites si l'Opération GAROUA l'a atteinte, et par quel événement.
\end{boxexo}

\begin{boxexo}[title={\ding{45}\quad Exercice~P.4. Le redémarrage \niveaumoyen}]
Expliquez à M.~ABENA, en cinq lignes, pourquoi il n'aurait pas dû
redémarrer les serveurs, et ce qu'il aurait dû faire à 22~h~05.
\end{boxexo}

\begin{boxexo}[title={\ding{45}\quad Exercice~P.5. Le prestataire \niveaumoyen}]
Listez trois clauses qui auraient dû figurer dans le contrat de SAHEL
INFOGÉRANCE, et dites pour chacune quel maillon de l'attaque elle
touchait.
\end{boxexo}

\begin{boxexo}[title={\ding{45}\quad Exercice~P.6. Les sauvegardes \niveaumoyen}]
Pourquoi des sauvegardes stockées sur un partage du domaine ne protègent
pas contre un rançongiciel~? Proposez une organisation qui y résiste.
\end{boxexo}

\begin{boxexo}[title={\ding{45}\quad Exercice~P.7. Ce que l'on ignore \niveaumoyen}]
Le lundi matin, rédigez la liste des questions auxquelles LiZbiZ répond
\NEK{}, et pour chacune la trace qui permettrait d'y répondre.
\end{boxexo}

\begin{boxexo}[title={\ding{45}\quad Exercice~P.8. Le temps de l'attaquant \niveaumoyen}]
Calculez la durée entre la première connexion de l'attaquant et le
chiffrement. Qu'est-ce que cette durée dit des capacités de détection de
LiZbiZ~?
\end{boxexo}

\begin{boxexo}[title={\ding{45}\quad Exercice~P.9. La note au directeur \niveauexpert}]
Rédigez la note d'une page que M.~KOUNOU remet à M.~ETEME le lundi à
8~h~: ce que l'on sait, ce que l'on ignore, les trois décisions à prendre
et, pour chacune, ce qu'il recommande.
\end{boxexo}

\begin{boxexo}[title={\ding{45}\quad Exercice~P.10. Le premier plan \niveauexpert}]
Proposez à LiZbiZ les cinq premières mesures à prendre dans le mois qui
suit, classées par priorité, et justifiez l'ordre.
\end{boxexo}

\bigskip

\begin{boxinfo}
\textbf{Pour aller plus loin.}
\begin{itemize}
    \item Le cas LiZbiZ vu par l'auditeur~: UE INF3226, Prologue et
          Chapitre~5.
    \item L'article sur l'instrument \OUI{}~/ \NON{}~/ \NEK{}~\cite{minka2026nek}.
\end{itemize}
\end{boxinfo}

\begin{boxinfo}
\textbf{Transition.} Avant de protéger, il faut savoir ce que l'on protège
et contre quoi. La Partie~\ref{part:fondements} pose les fondements~: les
propriétés de sécurité, les menaces et les attaques, la cryptographie, et
les leçons des grandes attaques.
\end{boxinfo}
